\documentclass{beamer}

\usepackage[T1]{fontenc}
\usepackage[utf8]{inputenc}
\usepackage[english]{babel}
\usepackage{lmodern}
\usepackage{booktabs}
\usepackage{multirow}
\usepackage{graphicx}
\usepackage{amsmath}

% Unipd theme (see main.tex for the available options)
\usetheme[pageofpages=of]{Unipd}

% Author and date are printed on the red title page: force them to white
\setbeamercolor{author}{fg=white}
\setbeamercolor{date}{fg=white}

% Figures: local copies of the manuscript ones, in ./Imm
\graphicspath{%
    {./}%
    {./Imm/}%
}

\title[Interface circuits for a six-phase VREH]{Interface Circuits for a Six-Phase\\Variable Reluctance Energy Harvester}
\subtitle{A comparative study of rectification and power conditioning topologies}
\author[Jacopo Garau]{Jacopo Garau\\[0.8ex]
    {\small Supervisor: Prof.~Alessandro Pozzebon}\\
    {\small Co-supervisors: Prof.~Sebastian Bader, Dr.~Ye Xu}}
\date{\small Master Thesis in Electronic Engineering --- A.Y.~2025/2026}

\begin{document}

% ---------------------------------------------------------------- 1. Title
\frame{\titlepage}

\section{Motivation and Objectives}

% ---------------------------------------------------------------- 2
\begin{frame}{Why energy harvesting?}
    \begin{itemize}
        \item Wireless sensor networks and IoT devices have grown very fast, and are
              widely used to \alert{collect data} that can then be processed for
              many different purposes.
        \item \alert{Millions of small nodes} therefore need to be powered.
        \item Batteries are expensive to maintain, have a finite life, and in many
              embedded or rotating installations \alert{cannot be replaced at all};
              they are also a poor choice for the environment.
    \end{itemize}
    \vfill
    \begin{block}{Energy harvesting}
        Convert energy already available in the environment into electrical power,
        so that the node powers itself for its whole lifetime.
    \end{block}
    \vfill
    This thesis targets \textbf{rotating motion}, converted through
    \textbf{Variable Reluctance Energy Harvesting (VREH)}.
\end{frame}

% ---------------------------------------------------------------- 3
\begin{frame}{The variable reluctance principle}
    \begin{columns}[T]
        \column{0.46\textwidth}
            \small
            \begin{itemize}
                \item Magnet and coil stay \alert{completely still}.
                \item Only a toothed wheel rotates across the air gap.
                \item The passing teeth modulate the reluctance, hence the flux
                      linked by the coil: an EMF is induced.
            \end{itemize}
            \vspace{0.5ex}
            \textbf{Advantages}
            \begin{itemize}
                \item No moving electrical parts.
                \item Easily scaled to different shaft sizes.
                \item Very small counter torque.
            \end{itemize}
        \column{0.52\textwidth}
            \centering
            \includegraphics[width=\textwidth]{VR_theory.pdf}\\[1ex]
            {\scriptsize Stationary pickup unit (magnet, pole piece, coil)
             facing the rotating toothed wheel.}
    \end{columns}
\end{frame}

% ---------------------------------------------------------------- 4
\begin{frame}{The six-phase MP-VREH}
    From single-unit prototypes for railway monitoring to an optimised
    \alert{six-phase} design for large vehicles (buses): six pickup units,
    electrically shifted by $60^\circ$, mounted around the shaft.
    \vfill
    \begin{table}
        \centering
        \small
        \begin{tabular}{cccccc}
            \toprule
            rpm & $f$ (Hz) & $V_\text{Lpk}$ (V) & $R_\text{coil}$ ($\Omega$) & $L_\text{coil}$ (mH) & $P_o$ (W)\\
            \midrule
            100 & 36.7  & 0.91 & 6 & 21.68 & 0.41\\
            200 & 73.3  & 1.35 & 6 & 21.46 & 0.91\\
            300 & 110.0 & 1.59 & 6 & 21.27 & 1.26\\
            400 & 146.7 & 1.74 & 6 & 21.12 & 1.51\\
            \bottomrule
        \end{tabular}
    \end{table}
    \vfill
    Target operating range: \alert{100--400 rpm}, i.e.\ 20--80 km/h for a typical
    large commercial vehicle. Each phase delivers less than $2\ \text{V}$ peak:
    the rectifier choice becomes critical.
\end{frame}

% ---------------------------------------------------------------- 5
\begin{frame}{Research gap and question}
    \begin{itemize}
        \item More phases mean more power, but a much harder electronics problem:
              \alert{six AC signals}, all shifted in time, all to be rectified.
        \item Previous work compared rectifier topologies for a
              \alert{single-phase} VREH only.
        \item How to condition the power of a six-phase harvester has never been
              studied systematically.
    \end{itemize}
    \vfill
    \begin{alertblock}{Research question}
        Which interface circuit extracts the most power from the six-phase
        MP-VREH, across its entire 100--400 rpm speed range?
    \end{alertblock}
\end{frame}

% ---------------------------------------------------------------- 6
\begin{frame}{Objectives and approach}
    A complete path, from theory to a working prototype:
    \vfill
    \begin{enumerate}
        \item \textbf{Model} the harvester: electrical model of one coil,
              extended to the six-phase system.
        \item \textbf{Propose} rectifier topologies able to combine six phases.
        \item \textbf{Simulate} each candidate in MATLAB/Simulink, sweeping the
              compensation capacitor and the load.
        \item \textbf{Design and build} a single reconfigurable prototype board.
        \item \textbf{Measure} every topology on the same hardware, on the same
              test bench.
        \item \textbf{Assess} compatibility with commercial power management ICs.
    \end{enumerate}
\end{frame}

\section{Modelling and Proposed Topologies}

% ---------------------------------------------------------------- 7
\begin{frame}{Coil model and matching}
    \begin{columns}[T]
        \column{0.52\textwidth}
            \small
            \begin{itemize}
                \item Each coil: ideal source in series with $R_\text{Coil}$ and
                      $L_\text{Coil}$.
                \item At low speed a purely \alert{resistive} load is enough to
                      extract the maximum power.
                \item As speed (and frequency) grows, $\omega_e L_\text{Coil}$
                      can no longer be neglected: the matched-load condition is lost.
                \item A \alert{series compensation capacitor} $C_m$ cancels the
                      inductive reactance and restores it:
            \end{itemize}
            \begin{equation*}
                C_m = \frac{1}{\omega_e^{2}\,L_\text{Coil}}
            \end{equation*}
            valid at one frequency at a time.
        \column{0.46\textwidth}
            \centering
            \includegraphics[width=0.95\textwidth]{Coil_model.pdf}\\[2ex]
            \includegraphics[width=0.85\textwidth]{Matching_condition.pdf}
    \end{columns}
\end{frame}

% ---------------------------------------------------------------- 8
\begin{frame}{The six real phase voltages}
    \begin{columns}[T]
        \column{0.58\textwidth}
            \centering
            \includegraphics[width=\textwidth]{OCV_100_focus.pdf}
        \column{0.40\textwidth}
            \begin{itemize}
                \item Measured open-circuit voltage of the six coils at
                      $100\ \text{rpm}$.
                \item The $60^\circ$ electrical shift between phases is clearly
                      visible.
                \item The waveform is \alert{not} a clean sine wave, and the six
                      units are not perfectly identical.
            \end{itemize}
    \end{columns}
\end{frame}

% ---------------------------------------------------------------- 9
\begin{frame}{From waveform to simulation}
    \begin{columns}[T]
        \column{0.44\textwidth}
            \begin{itemize}
                \item The RMS voltage of each coil was computed as a function of
                      speed.
                \item Each source was then replaced by an \alert{equivalent
                      sinusoid}, scaled through the measured crest factor.
                \item It delivers the same power as the real waveform, so the
                      simulation stays realistic while remaining tractable.
            \end{itemize}
        \column{0.54\textwidth}
            \centering
            \includegraphics[width=\textwidth]{Voltage_comp.pdf}\\[0.5ex]
            {\scriptsize FEM prediction, raw measured peak and crest-factor
             corrected value.}
    \end{columns}
\end{frame}

% ---------------------------------------------------------------- 10
\begin{frame}{Strategy 1: FWR bridges}
    \begin{columns}[T]
        \column{0.52\textwidth}
            \small
            The six phases split into \alert{two three-phase groups}
            ($k=0,2,4$ and $k=1,3,5$), each feeding a standard three-phase bridge.
            \vspace{1ex}
            \begin{itemize}
                \item \textbf{Star (Y)}: line voltage is $\sqrt{3}$ times the
                      phase EMF, so the fixed two-diode drop costs relatively less.
                \item \textbf{Delta ($\Delta$)}: line voltage equals the coil EMF,
                      no amplitude gain, but no floating node.
            \end{itemize}
            \vspace{1ex}
            Expectation: \alert{higher open-circuit voltage for the star
            connection}. Both are carried forward as candidates.
        \column{0.46\textwidth}
            \centering
            \includegraphics[width=0.92\textwidth]{FWR.pdf}\\[1ex]
            {\scriptsize Three-phase bridge for one group; the current path always
             crosses two diodes.}
    \end{columns}
\end{frame}

% ---------------------------------------------------------------- 11
\begin{frame}{Strategy 2: NVC bridges}
    \begin{columns}[T]
        \column{0.52\textwidth}
            \footnotesize
            The six phases split into \alert{three two-phase groups}, each pairing
            phases $180^\circ$ apart ($k=0,3$; $1,4$; $2,5$) in \textbf{anti-series}:
            \begin{equation*}
                v_A - v_B = e_k - e_{k+3} = 2\,e_k
            \end{equation*}
            ideally double the voltage, at the cost of a doubled series impedance.
            \vspace{0.5ex}
            \begin{itemize}
                \item Four MOSFETs with \alert{cross-coupled gates}: two conduct at
                      a time, no gate driver needed.
                \item They bypass the parallel Schottky diodes, replacing the fixed
                      drop with a small $R_\text{on}$; below
                      $V_\text{off}\approx1.2\ \text{V}$ the bridge falls back to
                      diode behaviour.
            \end{itemize}
        \column{0.46\textwidth}
            \centering
            \includegraphics[width=0.95\textwidth]{NVC.pdf}\\[1ex]
            {\scriptsize Negative voltage converter for one anti-series pair.}
    \end{columns}
\end{frame}

\section{Prototype and Measurement Setup}

% ---------------------------------------------------------------- 12
\begin{frame}{The prototype board}
    \begin{columns}[T]
        \column{0.48\textwidth}
            \begin{itemize}
                \item Custom PCB designed in Altium Designer.
                \item \alert{Every configuration} (star, delta, NVC) selected by
                      moving jumpers, on the same hardware.
                \item Jumper-selectable capacitor bank to set $C_m$ of each coil.
                \item Assembled entirely by hand, with soldering iron and hot plate.
            \end{itemize}
        \column{0.50\textwidth}
            \centering
            \includegraphics[width=0.9\textwidth]{Top3D.png}\\[1ex]
            \includegraphics[width=0.72\textwidth]{Soldering.jpg}\\[0.5ex]
            {\scriptsize One NVC bridge during hand assembly.}
    \end{columns}
\end{frame}

% ---------------------------------------------------------------- 13
\begin{frame}{Test bench}
    \begin{columns}[T]
        \column{0.55\textwidth}
            \centering
            \includegraphics[width=\textwidth]{Setup.jpg}\\[0.5ex]
            {\scriptsize Induction motor (left) coupled to the MP-VREH (right).}
        \column{0.43\textwidth}
            \begin{itemize}
                \item Induction motor driven by an inverter: rotation speed set
                      precisely over the whole 100--400 rpm range.
                \item Voltages acquired with a Saleae differential acquisition
                      device.
                \item Load set by an eight-decade programmable resistor board.
            \end{itemize}
    \end{columns}
\end{frame}

% ---------------------------------------------------------------- 14
\begin{frame}{Current sensing}
    \begin{columns}[T]
        \column{0.46\textwidth}
            \footnotesize
            \begin{itemize}
                \item Six custom sensing boards (INA114AU) measure the drop across
                      a shunt resistor differentially.
                \item Every shunt characterised with four-wire measurements.
                \item \alert{Each sensor calibrated individually} against precision
                      instruments, into its own lookup table.
                \item This makes power figures trustworthy down to the
                      \alert{milliwatt} level.
            \end{itemize}
        \column{0.52\textwidth}
            \centering
            \includegraphics[width=0.45\textwidth]{Curr_sens.jpg}\\[0.8ex]
            \includegraphics[width=0.82\textwidth]{Curr_sens.pdf}\\[0.3ex]
            {\scriptsize Calibration curves of the six sensors.}
    \end{columns}
\end{frame}

\section{Results}

% ---------------------------------------------------------------- 15
\begin{frame}{Simulation results}
    \begin{columns}[c]
        \column{0.33\textwidth}
            \centering
            \includegraphics[width=\textwidth]{Y_Power_sweep.pdf}\\
            {\scriptsize Star}
        \column{0.33\textwidth}
            \centering
            \includegraphics[width=\textwidth]{D_Power_sweep.pdf}\\
            {\scriptsize Delta}
        \column{0.33\textwidth}
            \centering
            \includegraphics[width=\textwidth]{NVC_Power_sweep.pdf}\\
            {\scriptsize NVC}
    \end{columns}
    \vfill
    \begin{itemize}
        \item Output power across the speed range, for several fixed values of $C_m$.
        \item At $100\ \text{rpm}$, the \alert{worst case} that sets the lower bound
              of what the system can deliver, the NVC gives $69.0\ \text{mW}$:
              about $6\times$ the delta connection ($10.7\ \text{mW}$) and nearly
              double the star one ($36.7\ \text{mW}$).
        \item The NVC is also the most efficient topology at every speed.
    \end{itemize}
\end{frame}

% ---------------------------------------------------------------- 16
\begin{frame}{Measured power}
    \begin{columns}[c]
        \column{0.33\textwidth}
            \centering
            \includegraphics[width=\textwidth]{Y_Fs_tot.pdf}\\
            {\scriptsize Star}
        \column{0.33\textwidth}
            \centering
            \includegraphics[width=\textwidth]{D_Fs_tot.pdf}\\
            {\scriptsize Delta}
        \column{0.33\textwidth}
            \centering
            \includegraphics[width=\textwidth]{NVC_Fs_tot.pdf}\\
            {\scriptsize NVC}
    \end{columns}
    \vfill
    \begin{itemize}
        \item At $100\ \text{rpm}$ the NVC still wins: $57.7\ \text{mW}$, against
              $47.6$ (star) and $14.9\ \text{mW}$ (delta).
        \item Above roughly $250\ \text{rpm}$ the \alert{star connection overtakes
              the NVC}: $433.3$ against $425.7\ \text{mW}$ at $400\ \text{rpm}$.
        \item Simulation alone did not predict this crossover.
    \end{itemize}
\end{frame}

% ---------------------------------------------------------------- 17
\begin{frame}{Measured efficiency}
    \begin{columns}[c]
        \column{0.33\textwidth}
            \centering
            \includegraphics[width=\textwidth]{Y_Eff_1.pdf}\\
            {\scriptsize Star}
        \column{0.33\textwidth}
            \centering
            \includegraphics[width=\textwidth]{D_Eff_2.pdf}\\
            {\scriptsize Delta}
        \column{0.33\textwidth}
            \centering
            \includegraphics[width=\textwidth]{NVC_Eff_tot.pdf}\\
            {\scriptsize NVC}
    \end{columns}
    \vspace{1ex}
    \small
    \begin{itemize}
        \item The NVC has the \alert{best efficiency at every speed}, up to
              $\eta_\text{Rec}\approx0.70$ at $400\ \text{rpm}$.
        \item The star connection reaches $\approx0.55$--$0.57$, the delta
              connection only $\approx0.48$--$0.51$: the weakest performer, as
              expected from its lower line voltage.
        \item So the two rankings disagree: \alert{efficiency says NVC, power says
              star at high speed}.
    \end{itemize}
\end{frame}

% ---------------------------------------------------------------- 18
\begin{frame}{Simulation vs measurement}
    Averaged over the whole speed range, the measured power differs from the
    simulated one by \alert{about 8\% for the star connection, up to almost 30\%
    for the NVC}.
    \vfill
    \begin{table}
        \centering
        \small
        \begin{tabular}{lccc}
            \toprule
            & Star & Delta & NVC\\
            \midrule
            Mean deviation of $P_O$ from simulation & 8\% & 16\% & \alert{30\%}\\
            \bottomrule
        \end{tabular}
    \end{table}
    \vfill
    \begin{itemize}
        \item The ceramic compensation capacitors \alert{lose part of their
              capacitance} under the real operating voltage, which grows with speed:
              a larger nominal $C_m$ is needed to reach the same effective value.
        \item Tolerances of the real components, and the parasitic resistance of
              the bridges, add losses the ideal model does not include.
        \item The sources are not perfectly sinusoidal.
    \end{itemize}
\end{frame}

% ---------------------------------------------------------------- 19
\begin{frame}{PMIC compatibility}
    \footnotesize
    A PMIC conditions the rectified power and delivers a \alert{stable supply} to
    the load: for example a node whose sensors detect and prevent malfunctions of
    the system or the vehicle it is installed on.
    \begin{table}
        \centering
        \footnotesize
        \begin{tabular}{lccc}
            \toprule
            PMIC (configuration) & $V_{in}$ range (V) & $V_{cs}$ (V) & Star connection\\
            \midrule
            DS-AEM13920\_VR      & $0.120$--$4.6$ & $0.275$        & up to $250\ \text{rpm}$\\
            DS-AEM10941\_VR      & $0.05$--$5$    & $0.38$         & up to $250\ \text{rpm}$\\
            bq25570              & $0.1$--$5.1$   & $0.6$          & up to $250\ \text{rpm}$\\
            SPV1050 (Boost)      & $0.15$--$5.3$  & $0.55$--$0.58$ & up to $250\ \text{rpm}$\\
            SPV1050 (Buck-Boost) & $0.15$--$18$   & $2.6$--$2.8$   & full range\\
            SPV1040              & $0.3$--$5.5$   & $0.8$--$2$     & up to $300\ \text{rpm}$\\
            \bottomrule
        \end{tabular}
    \end{table}
    \begin{itemize}
        \item Delta and NVC stay below about $4\ \text{V}$ over the whole range:
              \alert{compatible with every configuration tested}.
        \item The star connection climbs to $7\ \text{V}$ at $400\ \text{rpm}$: it
              needs a wide-input converter or an input clamp.
    \end{itemize}
\end{frame}

\section{Conclusions}

% ---------------------------------------------------------------- 20
\begin{frame}{Conclusions and future work}
    \small
    \begin{block}{Conclusions}
        The NVC and the star connection are the two strongest candidates, each
        suited to a different part of the speed range; the delta connection is
        consistently the weakest. NVC and delta connect safely to every PMIC
        tested, while the star connection needs extra protection at high speed:
        its higher open-circuit voltage can damage the input of the PMIC.
    \end{block}
    \begin{alertblock}{Limitation}
        The harvester prototype is 3D-printed, and therefore sensitive to
        vibration, especially at high speed: this limits how repeatable the
        measurements can be.
    \end{alertblock}
    \begin{block}{Future work}
        An adaptive compensation network tracking the optimal $C_m$ as speed
        changes, and a full hardware test of the harvester connected to a real
        PMIC, to measure the true efficiency from the harvester to the final load.
    \end{block}
\end{frame}

% ---------------------------------------------------------------- 21
\begin{emptyframe}
    Thank you for your attention!
\end{emptyframe}

\appendix

\begin{frame}{Backup: star connection}
    \begin{table}
        \centering
        \small
        \begin{tabular}{cccccc}
            \toprule
            \multirow{2}{*}{rpm} & \multicolumn{2}{c}{$P_O$ (mW)} & & \multicolumn{2}{c}{$\eta_\text{Rec}$}\\
            \cmidrule{2-3}\cmidrule{5-6}
            & Simulated & Measured & & Simulated & Measured\\
            \midrule
            100 & 36.7  & 47.6  & & 0.430 & 0.423\\
            150 & 94.3  & 107.5 & & 0.541 & 0.475\\
            200 & 165.8 & 175.6 & & 0.608 & 0.507\\
            250 & 240.0 & 245.1 & & 0.650 & 0.529\\
            300 & 308.8 & 312.7 & & 0.674 & 0.544\\
            350 & 372.3 & 374.6 & & 0.691 & 0.556\\
            400 & 430.0 & 433.3 & & 0.710 & 0.565\\
            \bottomrule
        \end{tabular}
    \end{table}
\end{frame}

\begin{frame}{Backup: delta and NVC}
    \begin{columns}[T]
        \column{0.49\textwidth}
            \centering {\small Delta}\\[0.5ex]
            \scriptsize\setlength{\tabcolsep}{3pt}
            \begin{tabular}{ccccc}
                \toprule
                \multirow{2}{*}{rpm} & \multicolumn{2}{c}{$P_O$ (mW)} & \multicolumn{2}{c}{$\eta_\text{Rec}$}\\
                \cmidrule{2-3}\cmidrule{4-5}
                & Sim. & Meas. & Sim. & Meas.\\
                \midrule
                100 & 10.7  & 14.9  & 0.210 & 0.279\\
                150 & 43.7  & 43.2  & 0.340 & 0.379\\
                200 & 95.0  & 84.2  & 0.425 & 0.441\\
                250 & 158.3 & 134.1 & 0.482 & 0.480\\
                300 & 228.5 & 190.8 & 0.522 & 0.508\\
                350 & 298.0 & 248.7 & 0.550 & 0.515\\
                400 & 365.0 & 310.7 & 0.570 & 0.511\\
                \bottomrule
            \end{tabular}
        \column{0.49\textwidth}
            \centering {\small NVC}\\[0.5ex]
            \scriptsize\setlength{\tabcolsep}{3pt}
            \begin{tabular}{ccccc}
                \toprule
                \multirow{2}{*}{rpm} & \multicolumn{2}{c}{$P_O$ (mW)} & \multicolumn{2}{c}{$\eta_\text{Rec}$}\\
                \cmidrule{2-3}\cmidrule{4-5}
                & Sim. & Meas. & Sim. & Meas.\\
                \midrule
                100 & 69.0  & 57.7  & 0.665 & 0.568\\
                150 & 146.5 & 109.4 & 0.768 & 0.619\\
                200 & 238.9 & 175.9 & 0.821 & 0.648\\
                250 & 374.8 & 245.2 & 0.841 & 0.668\\
                300 & 508.0 & 308.0 & 0.842 & 0.682\\
                350 & 576.8 & 367.7 & 0.834 & 0.693\\
                400 & 610.0 & 425.7 & 0.820 & 0.701\\
                \bottomrule
            \end{tabular}
    \end{columns}
\end{frame}

\begin{frame}{Backup: open-circuit voltage}
    \begin{table}
        \centering
        \small
        \begin{tabular}{cccccc}
            \toprule
            \multirow{2}{*}{rpm} & \multicolumn{2}{c}{Star} & \multicolumn{2}{c}{Delta} & NVC\\
            \cmidrule{2-3}\cmidrule{4-5}\cmidrule{6-6}
            & $V_{OC}$ (V) & ripple & $V_{OC}$ (V) & ripple & $V_{OC}$ (V)\\
            \midrule
            100 & 1.98 & 0.012 & 1.000 & 0.022 & 2.70\\
            150 & 2.83 & 0.012 & 1.460 & 0.021 & 3.05\\
            200 & 3.67 & 0.011 & 1.910 & 0.017 & 3.20\\
            250 & 4.51 & 0.012 & 2.365 & 0.017 & 3.15\\
            300 & 5.32 & 0.012 & 2.825 & 0.017 & 3.15\\
            350 & 6.17 & 0.011 & 3.275 & 0.017 & 3.15\\
            400 & 7.04 & 0.012 & 3.755 & 0.017 & 3.15\\
            \bottomrule
        \end{tabular}
    \end{table}
    The NVC saturates at $3.15\ \text{V}$ above $200\ \text{rpm}$, with a ripple an
    order of magnitude larger than the FWR connections, consistent with its active
    switching operation.
\end{frame}

\end{document}
